Using $\sin2x=2\sin x\cos x$, the equation becomes $\sqrt x\sin x(\sqrt x\cos x-1)=0$. Thus all nonnegative multiples $x=k\pi$ are solutions, together with solutions of $\sqrt x\cos x=1$. There are no such extra roots in $[0 , \dfrac{\pi}{2}]$: for $0\le x\le1$, the product is less than 1, while for $1 \le x \le \dfrac{\pi}{2}$, $\sqrt{x} \cos x \le \sqrt{\dfrac{\pi}{2}} \cos 1 < 1$. In each later positive-cosine interval $((2 k - \dfrac{1}{2}) \pi , (2 k + \dfrac{1}{2}) \pi)$, $k\ge1$, there are exactly two extra roots, one on each side of the unique maximum. Indeed the derivative has the sign of $\cos x-2x\sin x$; it is positive on the left half and changes once on the right half because $\tan x = \dfrac{1}{2x}$ has exactly one root there. The function is zero at both interval endpoints and exceeds 1 at $2k\pi$. This fully characterises the additional transcendental roots.
